% Generated draft: requires mathematical and source review.
\documentclass[11pt]{article}
\usepackage[margin=25mm]{geometry}
\usepackage{fontspec}
\setmainfont{Latin Modern Roman}
\usepackage{amsmath,amssymb,mathtools,unicode-math}
\setmathfont{Latin Modern Math}
\usepackage{enumitem,needspace,xurl,hyperref}
\setlength{\parindent}{0pt}
\setlength{\parskip}{5pt}
\setlength{\emergencystretch}{3em}
\newcommand{\N}{\mathbb{N}}
\newcommand{\Z}{\mathbb{Z}}
\newcommand{\Q}{\mathbb{Q}}
\newcommand{\R}{\mathbb{R}}
\newcommand{\C}{\mathbb{C}}
\newcommand{\F}{\mathbb{F}}
\newcommand{\A}{\mathbb{A}}
\newcommand{\PP}{\mathbb{P}}
\newcommand{\E}{\mathbb{E}}
\newcommand{\eps}{\varepsilon}
\newcommand{\dd}{\,\mathrm{d}}
\begin{document}
\section*{UnsolvedMath ID: 10400171}
\textbf{Source problem code:} catalogue entry
\subsection{Definitions and mathematical statement}
Let \(\mathcal C_1\) be a semisimple ribbon category with degenerate \(S\)-matrix, satisfying the additional hypotheses of the source's Problem 9.6. In the usual finite setting the entries of \(S\) are categorical traces of double braidings between simple objects; degeneracy means that \(S\) is not invertible. Let \(\mathcal C_2\) be the modular tensor category obtained from \(\mathcal C_1\) by the particular method in [288]. Write \((V_1,Z_1)\) and \((V_2,Z_2)\) for the two TQFTs referred to in the source as arising from these categories.

Kawahigashi asks to determine their relationship. A precise comparison must first identify the two constructions and their bordism/extra-structure domains, and then state the relation between their state spaces, cobordism maps and closed-manifold values. The source specifies no equality, natural transformation, quotient, extension or numerical formula, so none is asserted as its intended answer. In particular, one cannot use the ordinary modular-category construction to define the first theory merely by treating \(S_1\) as invertible: its degeneracy is an explicit hypothesis, and the missing Problem 9.6 must explain the first TQFT. Likewise the method [288] is part of the input, not an arbitrary choice among modular categories associated with \(\mathcal C_1\).
\subsection{Short English statement}
Compare the TQFT from a specified degenerate semisimple ribbon category with the TQFT from its modular category constructed by [288]. The missing hypotheses and construction references are essential to identify both theories.
\subsection{Sources}
\begin{itemize}
\item Local UnsolvedMath registry, permanent ID \(10400171\), exact source code \texttt{catalogue entry}.
\item Y. Kawahigashi, attribution in Ohtsuki, \emph{Problems on invariants of knots and 3-manifolds} (2002), Problem 9.7, PDF page 140: \url{https://msp.org/gtm/2002/04/gtm-2002-04-024s.pdf}. Recorded access 2026-07-29, pdf-font-structure extraction.
\item Problem 9.6 and internal reference [288] are invoked by the source; their content and full bibliography are not supplied.
\end{itemize}
\end{document}
